% mage-auto-heading: return-governability-and-authority
% mage-auto-heading: the-expanding-frontier-of-explicit-engineering
% mage-auto-heading: scope-beyond-software
% Authoritative LaTeX manuscript; Markdown is a generated compatibility projection.

% <!-- part-title: The Theory -->

% <!-- chapter-title: Scope Conditions -->

\chapter{Scope Conditions}

% <!-- point: scope-conditions-bound-where-theory-applies | MAGE requires a worthwhile investment, obligations that can be governed adequately and organizational capacity to act on diagnosis. -->

The four propositions are conditional. A representation can be expensive to keep current, a control can enforce the wrong obligation, and an organization can recognize a structural problem without being able to change it. These limits are part of the theory: they identify where the proposed engineering moves should earn their cost and where another arrangement may be better.

\section{Return, Governability, and Authority}\label{heading:return-governability-and-authority}

% <!-- point: mage-applicability-three-questions | Ask whether explicit engineering earns a return, whether the obligation can be governed adequately and whether the organization can change the shared environment. -->

Whether MAGE applies to a given engineering surface comes down to three questions:

\begin{enumerate}
\tightlist
\item \textbf{Will it earn a return?} Is there enough recurrence or consequence to repay the investment?
\item \textbf{Can the obligation be governed adequately?} Can the property be represented, evidenced, and evaluated well enough?
\item \textbf{Can the organization act on what it learns?} Can it distinguish structural pressure from local defects, and does it have authority to change the shared environment?
\end{enumerate}

% <!-- point: return-from-recurrence-or-consequence | Return comes usually from recurrence but also from consequence: does avoided future cost justify construction and carrying cost? -->

\textbf{Does the investment earn a return?} Recurrence is the usual source of return. If the same architectural knowledge, review judgment, or failure class will be encountered hundreds of times, making it explicit once can avoid paying for the same reasoning repeatedly. But recurrence is not the only justification. A rare obligation with catastrophic consequences may warrant extensive modeling and assurance because preventing a single failure can repay the investment. The economic question is therefore broader: \textbf{does the expected reduction in future reasoning, risk, or recovery cost justify the construction and carrying cost of the asset?}

% <!-- point: governability-is-explicit-evidenced-evaluable | Governability is how well an obligation can be made explicit, evidenced, and evaluated; deterministic checking is its strongest region, not its definition. -->

\textbf{Can the obligation be governed adequately?} Some properties can be represented precisely and settled mechanically: a schema conforms, a dependency is permitted, an invariant holds. Others admit statistical evidence but not a definitive predicate. Still others depend substantially on aesthetic, scientific, organizational, or social judgment. Governability is the degree to which an important obligation can be made explicit, supplied with relevant evidence, and evaluated with adequate reliability and acceptable cost. Deterministic checking is its strongest region, not its definition.

% <!-- point: residual-probabilistic-surface-is-not-necessarily-debt | The residual probabilistic surface is not necessarily debt; some judgment belongs there, so do not minimize it blindly. -->

The residual probabilistic surface is therefore not necessarily technical debt. Some judgment belongs there. The objective is not to minimize it blindly, but to avoid repeatedly purchasing expensive probabilistic reasoning or residual risk where a cheaper, adequate, and trustworthy representation or mechanism can do better.

% <!-- point: authority-connects-diagnosis-to-the-shared-environment | Diagnostic capacity and authority connect structural lessons to the shared environment; source access alone does not supply that authority. -->

The organization also needs diagnostic capacity and authority to adapt. Repeated failures become useful evidence only when someone can recognize the missing abstraction, boundary, invariant, or control. That diagnosis must then have a path into the shared environment. A contractor may identify a recurring failure without permission to modify the client's development pipeline. A packaged agent system may allow task instructions while preventing changes to its underlying controls. The engineer can see the problem in either case, yet be unable to convert the diagnosis into a durable intervention.

MAGE therefore requires sufficient control over the engineering environment rather than ownership of the software being developed. An engineer working on a proprietary platform may hold extensive governance authority, while a contractor with full source-code access may hold almost none. Authority need not reside in one person; it has to connect diagnosis to the environment.

% <!-- point: most-systems-mix-surfaces-which-deserve-enforcement | Most systems mix surfaces; the question is which properties deserve durable representation or enforcement, not whether a project uses MAGE. -->

A single system can contain obligations suited to durable enforcement alongside obligations that still require expert judgment. The practical question is therefore not whether a project ``uses MAGE,'' but which properties deserve durable representation or enforcement, how strong the evidence should be, and which decisions should remain with experts.

\section{The Expanding Frontier of Explicit Engineering}\label{heading:the-expanding-frontier-of-explicit-engineering}

% <!-- point: consequence-does-not-imply-governability | A consequential obligation can justify explicit engineering and careful judgment without becoming mechanically decidable. -->

Consequence does not imply governability. Some important obligations admit precise models and strong evaluation: structural loads, resource bounds, protocol states, or dependency relations may support calculation, exhaustive analysis, or proof. Other obligations remain substantially qualitative even when an organization has strong reason, or a legal duty, to evaluate them carefully.

% <!-- box-family: inset -->

% <!-- inset-domain: ENGINEERING PRACTICE -->

% mage-box-prefix: Inset —

\begin{bookbox}{inset}{Hard Obligations, Judgment Calls}
\label{inset-hard-obligations-judgment-calls}
An explicit, consequential obligation can still require judgment. Consider a historic building. The Secretary of the Interior\textquotesingle s Standards for Rehabilitation require certified projects to preserve historic character, and new work must be compatible with the historic building in its massing, size, scale, and architectural features. Replacement features may need to match the original design, color, texture, and other visual qualities.\citep{nps-rehab-standards} These are real requirements that can determine whether a project is approved, but there is no complete test for ``compatible with the historic character.'' Engineers and architects can measure dimensions, identify materials, record features that must be preserved, compare photographs and drawings, and use precedent and design rules. These make the judgment better informed and more consistent. They do not remove the judgment.

Software has the same problem. The Department of Justice\textquotesingle s 2024 accessibility rule generally requires state and local governments to make their web and mobile content conform to \magevalue{doj_technical_standard}, including conventional electronic documents where the criteria apply.\citep{doj-web-accessibility-2024} Many resulting requirements can be checked mechanically, but the rule also permits an alternative technique when it provides ``substantially equivalent or greater accessibility and usability.'' Whether it does so can depend on the information and functionality preserved, the users and tasks involved, and the barriers the alternative introduces. DocAble sits on this boundary. It can mechanically check many properties that its document models make observable and use probabilistic models for semantic work such as interpreting figures and generating descriptions. But neither proves that every blind student will receive a substantially equivalent experience from every remediated document. Some judgment remains.

The historic-building and accessibility examples make the same point. An obligation can be important, explicit, and legally binding without having a complete formal test. Engineering can still make the judgment better: make the obligation explicit, mechanize the parts that can be checked, and gather good evidence for the rest. The boundary can also move. Better models and measurements can turn judgments that once required an expert into properties the environment can check; better GenAI may make us willing to entrust increasingly difficult judgments to an artificial expert. Here, expert does not necessarily mean human: it means a reasoner capable of making the judgment well enough for the consequence at stake. The goal is therefore not to eliminate judgment, but to give each obligation the strongest form of evaluation and enforcement the available evidence can support.
\end{bookbox}

% <!-- point: value-and-governability-are-two-dimensions-of-investment | The value of explicit engineering and feasible governability are distinct; valuable judgment may warrant investment without a complete predicate. -->

The \textbf{value of explicit engineering} asks how much explicit structure the obligation warrants. Recurrence, consequence, and the reasoning or recovery it saves can all raise that value. \textbf{Feasible governability} asks what degree of evaluation and enforcement the available representation and evidence can support. Where both are high, formalization and enforcement may pay. Where value is high but governability is limited, investment can instead improve the models, evidence, procedures, and review through which judgment occurs. \autoref{return-governability} shows the two dimensions together.

\begin{figure}[htbp]
\centering
\includegraphics[width=\linewidth]{book/latex/assets/return-governability-5cf809fe48.pdf}
\caption{\emph{The economics of explicit engineering.} The value of making an obligation explicit determines how much investment it warrants; feasible governability determines how much of that investment can become enforcement. High-value obligations can justify substantial engineering even when they remain judgmental, while highly governable obligations need not be mechanized when little value would result.}
\label{return-governability}
\end{figure}

% <!-- point: ci-lowers-the-cost-of-explicit-structure | Where commodity intelligence lowers construction and maintenance cost without sacrificing quality or useful life, explicit engineering can repay its cost at lower recurrence or consequence. -->

Commodity intelligence can change where work falls in this design space in two ways. First, it may lower the cost of constructing and maintaining explicit structure. Where it does, models, skills, specifications, validators, proofs, and other engineering assets can repay their cost at lower levels of recurrence or consequence, provided their quality and useful life are preserved. More obligations may then justify explicit engineering even though the obligations themselves have not changed.

% <!-- point: ci-can-increase-governability | Commodity intelligence can also increase governability, turning expert judgment into narrower evaluable or decidable properties. -->

Second, better representation, instrumentation, and evaluation can increase governability. An obligation initially expressed only as expert judgment may yield narrower properties that can be evaluated reliably; a property reconstructed semantically from implementation may become a predicate once an appropriate model exposes it directly. \magevalue{chapter:one-problem-many-models} described the practical version: make obligations enforceable where feasible. The determinization frontier describes its strongest case: Modeling changes the representation until a previously reconstructed judgment becomes mechanically decidable, after which Alignment can carry the decision repeatably.

% <!-- point: movements-need-not-end-in-formalization | These movements need not end in formalization; explicit engineering still pays when it makes judgment better informed and repeatable. -->

These two changes need not end in the same place. Some questions resist enforcement because the necessary representation or evidence remains unavailable. Others admit probabilistic evaluation adequate for one consequence but not another. Still others are intrinsically plural or contextual: no amount of cheaper implementation turns good writing, scientific importance, or ``compatible with the historic character'' into a single definitive predicate. Explicit engineering still pays when it makes those judgments better represented, better informed, more repeatable, or easier to review.

% <!-- point: the-frontier-is-not-a-destination-evolution-is-an-outer-loop | The boundary moves without converging on full specification; implementation also discovers requirements, and evolution can reopen settled questions. -->

Neither movement implies a destination for software engineering as a whole. Many consequential questions remain open because the system's appropriate shape is still being discovered: which abstraction matters, which tradeoff users will accept, or which failure will dominate in operation. Implementation can be an experiment that produces information about the specification. Even a question that is closed today may reopen as users, dependencies, workloads, regulations, or adversaries change. Software evolution remains an outer loop around increasingly capable episodes of automated realization.\citep{lehman2003evolution}

The failure mode is \textbf{premature closure}. An explicit representation can make the wrong target easier to optimize. A small probabilistic surface is therefore not itself evidence of good engineering. MAGE asks which questions can be governed adequately now, which still need investigation, and what structure would help the next decision. It does not predict convergence toward complete specification.

\section{Scope Beyond Software}\label{heading:scope-beyond-software}

% <!-- point: scope-beyond-software-is-a-conditional-generalization | The mechanisms may transfer where useful representations, trustworthy correspondence and adequate governance are available; that transfer remains untested here. -->

The mechanisms described here do not intrinsically require software. Other engineering work also uses representations, independent evidence, reusable procedures, and inherited judgment. Transfer requires useful representations of the consequential state, credible correspondence between those representations and the system, and mechanisms adequate to govern the relevant obligations. Where those conditions are weak, better agents or expert judgment may remain the appropriate means of doing the work.

These scope conditions are theoretical, informed by \magevalue{chapter:software-factories}; they are not estimates of how often MAGE will succeed across a population. \magevalue{chapter:profession} examines other domains and what the account implies for engineers. The research agenda asks how to test the relationships first, before treating that wider application as established.
